\documentclass{article}
\usepackage{mathtools} 
\usepackage{fontspec}
\usepackage[UTF8]{ctex}
\usepackage{amsthm}
\usepackage{mdframed}
\usepackage{xcolor}
\usepackage{amssymb}
\usepackage{amsmath}


% 定义新的带灰色背景的说明环境 zremark
\newmdtheoremenv[
  backgroundcolor=gray!10,
  % 边框与背景一致，边框线会消失
  linecolor=gray!10
]{zremark}{说明}

% 通用矩阵命令: \flexmatrix{矩阵名}{元素符号}{行数}{列数}
\newcommand{\flexmatrix}[4]{
  \[
  #1 = \begin{pmatrix}
    #2_{11}     & #2_{12}     & \cdots & #2_{1#4}   \\
    #2_{21}     & #2_{22}     & \cdots & #2_{2#4}   \\
    \vdots      & \vdots      & \ddots & \vdots     \\
    #2_{#31}    & #2_{#32}    & \cdots & #2_{#3#4}
  \end{pmatrix}
  \]
}

% 简化版命令（默认矩阵名为A，元素符号为a）: \quickmatrix{行数}{列数}
\newcommand{\quickmatrix}[2]{\flexmatrix{A}{a}{#1}{#2}}

\begin{document}
\title{注释 6.3}
\author{张志聪}
\maketitle
% \begin{zremark}
%   注2

%   查了下资料，这里指的应该是: 泰勒展开的系数是唯一的。
% \end{zremark}

% \textbf{证明：}

% 反证法，假设不唯一，存在两个$p_n(x)$形式的多项式：
% \begin{align*}
%   f(x) & = a_0 + a_1(x - x_0) + a_2(x - x_0)^2 + \cdots + a_n(x - x_0)^n + o((x - x_0)^n) \\
%   f(x) & = b_0 + b_1(x - x_0) + b_2(x - x_0)^2 + \cdots + b_n(x - x_0)^n + o((x - x_0)^n)
% \end{align*}
% 两式相减
% \begin{align*}
%   0 = (a_0 - b_0) + (a_1 - b_1)(x - x_0) + (a_2 - b_2)(x - x_0)^2 + \cdots + (a_n - b_n)(x - x_0)^n
%   + o((x - x_0)^n)
% \end{align*}

\begin{zremark}
  无穷小量间的关系
\end{zremark}

\begin{itemize}
  \item $o(x^n) + o(x^n) = o(x^n) \ (x \to 0)$;
  \item $o(x^n) \cdot o(x^n) = o(x^n) \ (x \to 0)$;
  \item $x^k \cdot o(x^n) = o(x^{k + n}), k \in \mathbb{N^+} \ (x \to 0)$;
  \item $x^k \cdot o(x^n) = o(x^{n}), k \in \mathbb{N^+} \ (x \to 0)$;
  \item $x^{n + k} = o(x^n), k \in \mathbb{N^+} \ (x \to 0)$;
  \item $-1 \cdot o(x^n) = o(x^n) \ (x \to 0)$.
\end{itemize}


\end{document}